% BHU PYQ -- Manually transcribed & error-corrected
% Paper: MTB-202 : Statics and Dynamics
% Exam: B.A. (Hons.) Semester II Examination, 2022-23

\documentclass[11pt,a4paper]{article}
\usepackage[a4paper,top=2.5cm,bottom=2.5cm,left=2.8cm,right=2.8cm,headheight=14pt]{geometry}
\usepackage{amsmath,amssymb}
\usepackage[T1]{fontenc}
\usepackage[utf8]{inputenc}
\usepackage{lmodern}
\usepackage{microtype}
\usepackage{fancyhdr}
\usepackage{enumitem}
\usepackage{hyperref}
\usepackage{lastpage}
\usepackage{parskip}

\hypersetup{colorlinks=true,urlcolor=black,linkcolor=black,
  pdftitle={MTB-202: Statics and Dynamics -- BHU PYQ},pdfauthor={Banaras Hindu University}}

\pagestyle{fancy}\fancyhf{}
\renewcommand{\headrulewidth}{0.4pt}\renewcommand{\footrulewidth}{0.4pt}
\fancyhead[L]{\small   \textit{Banaras Hindu University}}
\fancyhead[R]{\small   \textit{MTB-202: Statics and Dynamics}}
\fancyfoot[C]{\small Page  \thepage\ of \pageref{LastPage}}


\newlist{parts}{enumerate}{2}
\setlist[parts,1]{label=(\alph*),leftmargin=*,itemsep=5pt,topsep=3pt}
\setlist[parts,2]{label=(\roman*),leftmargin=*,itemsep=3pt,topsep=2pt}

\newcommand{\pts}[1]{\hfill   \textit{[#1]}}


\begin{document}


\begin{center}
  {\large\bfseries BANARAS HINDU UNIVERSITY}\\[2pt]
  {
\normalsize B.A. (Hons.) --- Semester II Examination, 2022-23}\\[6pt]
  \rule{0.8\linewidth}{0.5pt}\\[6pt]
  {\large\bfseries Paper No. MTB-202: Statics and Dynamics}\\[4pt]
  {
\normalsize Time: 3 Hours \hfill Maximum Marks: 70}
\end{center}
\rule{\linewidth}{0.4pt}
\smallskip



\noindent   \textit{Instructions: Answer   \textbf{five} questions including Question~1 which is compulsory. Symbols have their usual meanings.}
\bigskip




\noindent   \textbf{Question 1.}   \textit{(Compulsory --- 3 marks each for (a)-(d), 2 marks each for (e)-(i).)}

\begin{parts}
  \item Determine the relation between angular and linear velocity for a particle moving in a curved path. \pts{3}
  \item Define moment of a force about a point. In a rigid body, the forces $\vec{P} = -2\vec{i} + 2\vec{j}$, $\vec{Q} = 3\vec{i} + 4\vec{j}$ and $\vec{R} = -3\vec{i} + \vec{j}$ act at the points $(1, 2)$, $(-1, -1)$ and $(1, -2)$, respectively. Determine whether the body will rotate or not. \pts{3}
  \item Three forces $P, Q$ and $R$ act along the sides $AB, BC, CA$ of a triangle $ABC$, respectively, formed by the coordinates $A(a, 0)$, $B(0, a)$ and $C(0, 0)$. Find the line of action of the resultant force. \pts{3}
  \item Show that the areal velocity is constant for a particle moving in a central orbit. \pts{3}
  \item An object is moving in a plane which moves with a uniform velocity. Justify whether or not the particle is moving in a Newtonian frame. \pts{2}
  \item State Kepler's laws of planetary motion. \pts{2}
  \item If a point moves along a circle, show that its angular velocity about any point on the circle is half of that about the centre. \pts{2}
  \item Define epoch and argument in a simple harmonic motion. \pts{2}
  \item A particle describes the parabola $p^2 = ar$ under a force directed towards the focus. Find the velocity at any point on the parabola. \pts{2}
\end{parts}

\medskip\rule{\linewidth}{0.2pt}

\medskip


\noindent   \textbf{Question 2.}

\begin{parts}
  \item Write down the equations of motion of a particle moving in a central orbit under a central force of magnitude $F$ and deduce the differential equation of the orbit in pedal form. \pts{7}
  \item In a simple harmonic motion, if $f$ be the acceleration and $v$ be the velocity at any instant of time, show that $f^2 T^2 + 4\pi^2 v^2$ is constant, where $T$ is the time period of the oscillation. \pts{5}
\end{parts}

\medskip\rule{\linewidth}{0.2pt}

\medskip


\noindent   \textbf{Question 3.}

\begin{parts}
  \item A particle describes an ellipse under a force $\mu/r^2$ towards the focus. If it was projected with a velocity $V$ from a point distant $R$ from the centre of force, show that the periodic time is:
  \[
  T = \frac{2\pi}{\sqrt{\mu}} \left( \frac{2}{R} - \frac{V^2}{\mu} \right)^{-3/2}
  \] \pts{6}
  \item A particle is moving along a straight line under an attractive force proportional to the distance from a fixed point on the line. If the particle starts from a position at a distance $a$ from the fixed point with velocity $V$, find the time of reaching the fixed point. \pts{6}
\end{parts}

\medskip\rule{\linewidth}{0.2pt}

\medskip


\noindent   \textbf{Question 4.}

\begin{parts}
  \item Show that a system of coplanar forces acting at different points of a rigid body can be reduced to a single force and a couple. Hence determine the conditions of equilibrium of the body. \pts{6}
  \item Forces of magnitudes 1, 2, 3, 4, 5 and 6 act in order along the sides of a regular hexagon. Show that the line of action of the resultant is parallel to the $x$-axis. \pts{6}
\end{parts}

\medskip\rule{\linewidth}{0.2pt}

\medskip


\noindent   \textbf{Question 5.}

\begin{parts}
  \item Derive the tangential and normal components of velocity and acceleration of a particle moving in a plane curve. \pts{6}
  \item A particle describes the curve $r = a \sec^2( \theta/2)$ such that the cross-radial velocity is constant. Show that the radial acceleration is constant. \pts{6}
\end{parts}

\medskip\rule{\linewidth}{0.2pt}

\medskip


\noindent   \textbf{Question 6.}

\begin{parts}
  \item State the principle of virtual work. A rhombus is formed by four uniform rods each of weight $W$ and length $l$ with smooth joints. It rests symmetrically with its upper sides in contact with two smooth pegs at a distance $2a$ apart. A weight $2W$ is hung at the lowest point. If the sides of the rhombus make an angle $ \theta$ with the vertical, show by the principle of virtual work that $\sin^3 \theta = \frac{3a}{l}$. \pts{6}
  \item Discuss the motion of a particle moving in a uniformly rotating frame with respect to a fixed frame. \pts{6}
\end{parts}

\medskip\rule{\linewidth}{0.2pt}

\medskip


\noindent   \textbf{Question 7.}

\begin{parts}
  \item A particle of mass $m$ slides down a smooth vertical cycloid $s = 4a\sin\psi$ from the cusp. Write down the equations of motion and find the time taken by the particle to reach the bottom of the cycloid. What will be the normal reaction at that point? \pts{6}
  \item A heavy particle of weight $mg$ is projected with velocity $V$ from the lowest point of a smooth vertical circle of radius $a$. Discuss the motion of the particle. Find the velocity of projection so that the particle will reach the highest point. \pts{6}
\end{parts}

\vfill
\rule{\linewidth}{0.4pt}

\begin{center}\small
  \href{https://bhu-exam.vercel.app/}{Click here to visit our website}\\[4pt]
  \href{https://me-aryan.vercel.app/}{Click here to visit our contributor}\\[4pt]
  \href{https://quashan-me.vercel.app/}{Click here to visit our contributor}
\end{center}

\end{document}
